\section{Singularities along a section over a divisor}
\label{sec:horizontal-center-constraints}

We now replace the isolated poles by a pole along the image of a holomorphic section of the fibration over a divisor of the base. The selected sections must then vanish to a certain order along that divisor. As for point poles, the normal variables in the fiber shift the integrability threshold only by a bounded amount, so this order differs from the order imposed by a divisorial pole on the base by a bounded number. Restriction sequences along the divisor show that each unit of vanishing order costs $O(m^{q-1})$ sections, and kernel convergence again reduces to the equilibrium theorem on the base.

For the restriction sequences we need a bound for the dimensions of section spaces that is uniform over a range of twists growing linearly with $m$.

\begin{lemma}\label{lem:linear-range-section-bound}\label{lem:finite-family-section-bound}
Let $M$ be a compact complex manifold of dimension $d\geq0$, and let $E,F_1,\ldots,F_N$ be holomorphic line bundles on $M$. For each $C_0>0$ there is a constant $C_1$ such that
\[
 h^0\Bigl(M,E^m\otimes\bigotimes_{i=1}^NF_i^{j_i}\Bigr)\leq C_1m^d
 \qquad(m\geq1,\ j_i\in\mathbb Z,\ |j_i|\leq C_0m).
\]
\end{lemma}
\begin{proof}
For $d=0$, the manifold $M$ is a finite set and the dimension is at most its number of points. Let $d\geq1$. Fix smooth Hermitian metrics on $E$ and the $F_i$, a smooth positive volume form $\mathrm{d}V$, and finitely many coordinate charts with holomorphic frames, together with smaller subcharts that still cover $M$. For $|j_i|\leq C_0m$, the local weight of the metric on $G=E^m\otimes\bigotimes_iF_i^{j_i}$ has second derivatives bounded by $Cm$. At a point $p$ of a smaller subchart, multiply the frame by the exponential of a holomorphic affine function to remove the constant and real-linear terms of this weight at $p$. The new weight vanishes at $p$ and is bounded by $Cm|z-p|^2$ on a fixed neighborhood. A ball of radius $\varepsilon m^{-1/2}$ about $p$ lies in the chart, with $\varepsilon>0$ independent of $p$, $m$ and the $j_i$, and on this ball the weight is bounded. The submean inequality for the holomorphic coefficient of a section $s$ on this ball gives
\[
 |s(p)|^2\leq Cm^d\int_M|s|^2\,\mathrm{d}V,
\]
where all constants depend only on the fixed charts, metrics, volume and $C_0$. Evaluation at $p$ has a one-dimensional target, so for an $L^2$-orthonormal basis $s_1,\ldots,s_D$ of $H^0(M,G)$ its squared norm is $\sum_{\ell=1}^D|s_\ell(p)|^2\leq Cm^d$. Integrating over $M$ gives $D\leq C\operatorname{Vol}(M)m^d$.
\end{proof}

Let $\pi\colon X\to B$ be a surjective holomorphic submersion of connected compact complex manifolds, with connected fibers of dimension $r$, and let $q=\dim B$. Fix a holomorphic line bundle $A$ on $B$ with a smooth Hermitian metric $h_A$, and put $\theta=c_1(A,h_A)$. Let $Z\subset B$ be a nonempty smooth divisor, and suppose that $\pi$ has a holomorphic section $\sigma\colon Z\to X$ over $Z$; put $Y=\sigma(Z)$. Fix $c>0$ and a smooth metric $h_Z$ on $\mathcal O_B(Z)$, with curvature $\Theta_Z$, such that $\alpha:=\theta-c\Theta_Z>0$. Let $s_Z$ be the canonical section of $\mathcal O_B(Z)$, and set $\psi=c\log|s_Z|_{h_Z}^2$.

\Needspace{11\baselineskip}
Fix a Hermitian form $\omega$ on $X$, and suppose that $\varphi\in\PSH(X,\omega)$ is locally bounded on $X\setminus Y$ and that, in submersion coordinates $(x,z',w_1,\ldots,w_r)$ near each point of $Y$,
\begin{equation}\label{eq:horizontal-center-singularity}
 \begin{gathered}
 \varphi=c\log\left(|x|^2+\sum_{j=1}^r|w_j|^2\right)+O(1),\\
 Z=\{x=0\},\qquad Y=\{x=w_1=\cdots=w_r=0\},
 \end{gathered}
\end{equation}
with a locally bounded remainder. Here $x$ and $z'$ are coordinates on the base, and the vector $z'$ of the $q-1$ tangential coordinates is absent when $q=1$.

\Needspace{23\baselineskip}
\begin{theorem}\label{thm:horizontal-center-descent}
The function $\psi$ is $\theta$-plurisubharmonic, with positive non-pluripolar mass
\[
 \int_B\mathrm{MA}_\theta(\psi)=\int_B\alpha^q>0.
\]
For every integer $m\geq1$ put $\ell_m=\lfloor mc\rfloor$ and $k_m=\max\{\ell_m-r,0\}$. Under holomorphic descent, the spaces selected by $\varphi$ on $X$ and by $\psi$ on $B$ are
\[
 \begin{split}
 V_m&:=H^0\bigl(X,\pi^*A^m\otimes\mathcal I(m\varphi)\bigr)
       =H^0\bigl(B,A^m(-k_m Z)\bigr),\\
 W_m&:=H^0\bigl(B,A^m\otimes\mathcal I(m\psi)\bigr)
       =H^0\bigl(B,A^m(-\ell_m Z)\bigr)\subseteq V_m.
 \end{split}
\]
There is a constant $C$, depending only on the fixed data, such that
\[
 0\leq\dim V_m-\dim W_m\leq Cm^{q-1}.
\]
For $q=1$ the difference is at most $r\,\#Z$ for every $m$.
\end{theorem}
\begin{proof}
By the Poincar\'e--Lelong formula, $\theta+\ddc\psi=\alpha+c[Z]\geq0$. The divisorial measure identity \eqref{eq:ordinary-divisor-measure}, applied with the residual form $\alpha>0$ and the bounded part zero, gives $\mathrm{MA}_\theta(\psi)=\alpha^q$, which proves the mass formula.

Every section of $\pi^*A^m$ descends by \cref{prop:partial-kernel-submersion}. In the $r+1$ normal variables $(x,w)$, \cref{lem:smooth-center-multiplier-ideal} gives
\[
 \mathcal I(m\varphi)=(x,w_1,\ldots,w_r)^{k_m}.
\]
The local coefficient $g(x,z')$ of a pulled-back section does not depend on $w$. Expanding $g$ in powers of $x$, we see that $g$ lies in this ideal exactly when the coefficients of $x^j$ vanish for $j<k_m$, that is, when $x^{k_m}$ divides $g$. This is an identity of germs in all tangential variables, so it holds at every point of $Z$. Outside $Y$ the function $\varphi$ is locally bounded. On the base, the same lemma in the single normal variable $x$ gives $\mathcal I(m\psi)=\mathcal O_B(-\ell_mZ)$. These facts prove the two identities and the inclusion.

For $k_m\leq j<\ell_m$, the restriction sequence along $Z$,
\[
 \begin{split}
 0&\longrightarrow A^m(-(j+1)Z)\longrightarrow A^m(-jZ)\\
  &\longrightarrow \bigl(A^m(-jZ)\bigr)|_Z\longrightarrow0,
 \end{split}
\]
shows that passing from vanishing order $j$ to order $j+1$ along $Z$ loses at most $h^0(Z,(A^m(-jZ))|_Z)$ dimensions. Apply \cref{lem:linear-range-section-bound} on $Z$ with $E=A|_Z$ and $F_1=\mathcal O_B(-Z)|_Z$; since $0\leq j\leq cm$, each loss is at most $Cm^{q-1}$, uniformly in $j$. There are at most $r$ steps, which proves the estimate. When $q=1$, the divisor $Z$ is a finite set and each restriction space has dimension $\#Z$, which gives the last assertion directly.
\end{proof}

We call $\psi$ the \emph{reference potential}, $W_m$ the \emph{reference spaces}, and $\dim V_m-\dim W_m$ the \emph{defect}; the same terms are used for the analogous comparisons below. The following example shows that the order $m^{q-1}$ in \cref{thm:horizontal-center-descent} is sharp.

\Needspace{12\baselineskip}
\begin{example}\label{ex:horizontal-center-sharp-error}
For all integers $q,r,k\geq1$ and every real $0<c<k$, the hypotheses above are satisfied on $X=\mathbb P^q\times\mathbb P^r$ with $B=\mathbb P^q$ and $A=\mathcal O_{\mathbb P^q}(k)$, for suitable choices made in the proof. For $q\geq2$ the defect is
\[
 \dim V_m-\dim W_m
   =\frac{r(k-c)^{q-1}}{(q-1)!}m^{q-1}+O(m^{q-2}),
\]
and for $q=1$ it equals $\min\{r,\lfloor mc\rfloor\}$ for every $m$.
\end{example}
\begin{proof}
Use homogeneous coordinates $[Z_0:\cdots:Z_q]$ on the base and $[W_0:\cdots:W_r]$ on the fiber. Let $Z=\{Z_0=0\}$ and $Y=Z\times\{[1:0:\cdots:0]\}$, and let $\omega_B,\omega_F$ be the Fubini--Study forms normalized to have top integrals one. Define
\[
 R_B=\frac{|Z_0|^2}{\sum_{i=0}^q|Z_i|^2},\qquad
 R_F=\frac{\sum_{j=1}^r|W_j|^2}{\sum_{j=0}^r|W_j|^2},\qquad
 \varphi=c\log(R_B+R_F),\quad\psi=c\log R_B.
\]
The function $\varphi$ is smooth off $Y$ and satisfies \eqref{eq:horizontal-center-singularity}. Over a common denominator, the numerator of $R_B+R_F$ is
\[
 \sum_{j=0}^r|Z_0W_j|^2+\sum_{i=0}^q\sum_{j=1}^r|Z_iW_j|^2,
\]
whose logarithm is plurisubharmonic in each affine chart. Hence $\ddc\varphi\geq-c(\omega_B+\omega_F)$. Choose an integer $b\geq c$ and set $\omega=k\omega_B+b\omega_F$. This is a Hermitian form, and $\omega+\ddc\varphi\geq(k-c)\omega_B+(b-c)\omega_F\geq0$. On the base, $\theta=k\omega_B$, $\Theta_Z=\omega_B$ and $\alpha=(k-c)\omega_B>0$.

Put $\ell_m=\lfloor mc\rfloor$. Divisibility of homogeneous polynomials by powers of $Z_0$ gives the exact dimensions
\[
 \dim W_m=\binom{km-\ell_m+q}{q},\qquad
 \dim V_m=\binom{km-\max\{\ell_m-r,0\}+q}{q}.
\]
For large $m$, the two residual degrees differ by $r$ and equal $(k-c)m+O(1)$. Subtracting the two binomial coefficients, which are polynomials of degree $q$ in the residual degrees, gives the asserted leading coefficient and remainder. For $q=1$ the difference is $\ell_m-\max\{\ell_m-r,0\}=\min\{r,\ell_m\}$ for every $m$.
\end{proof}

The bounded shift of divisorial order thus has negligible normalized codimension. To compare the kernels, use the volumes, metrics and continuous weights $v\circ\pi$ of \cref{sec:fiber-point-kernels}, together with the notation $\beta_{X,m,v\circ\pi,\varphi}^{(q)}$ and $\mathcal F_{X,m}^{(q)}$ introduced there, and put $E(v)=E_{\theta,B}[\psi]_{\mathcal I}(v)$.

\Needspace{12\baselineskip}
\begin{corollary}\label{cor:horizontal-center-kernel-limit}
Assume that $B$ is K\"ahler. Then $\varphi$ satisfies the conclusions of \cref{prop:descended-kernel-limit}, with $n=q+r$, and the total variation distance between $\beta_{X,m,v\circ\pi,\varphi}^{(q)}$ and the lift of the kernel measure selected by $\psi$ on $B$ is $O(m^{-1})$. In particular,
\[
 \beta_{X,m,v\circ\pi,\varphi}^{(q)}
    \rightharpoonup\mathcal L\bigl(\mathrm{MA}_\theta(E(v))\bigr)
 \qquad(v\in C(B)).
\]
\end{corollary}
\begin{proof}
By \cref{thm:horizontal-center-descent}, $\psi$ has positive non-pluripolar mass, and the reference spaces $W_m$ satisfy $W_m\subseteq V_m$ and $\dim V_m-\dim W_m=O(m^{q-1})$. Hence \cref{prop:descended-kernel-limit} applies.
\end{proof}

Thus the selecting potentials of \eqref{eq:horizontal-center-singularity} admit the approximation asked for in \cref{qu:fiber-selected-approximation}, whenever the curvature margin $\theta-c\Theta_Z$ is positive.

\subsection{A bounded divisorial band from fiber integrability}
\label{sec:fiber-integrability-band}

The comparison of \cref{thm:horizontal-center-descent} does not need the exact local form \eqref{eq:horizontal-center-singularity}. It suffices to bound the selecting potential from below by divisorial poles on all of $X$, and from above by an anisotropic model in a single chart over each divisor. The lower bound gives the inclusion of the reference spaces into the selected spaces. The upper bound forces vanishing along each divisor to an order that falls short of the reference order by a bounded amount: separate powers of the normal variables replace the shift by the fiber dimension with a weighted shift. The selecting potential may then be singular along a positive-dimensional subset of a fiber, and it need not descend to the base. Several divisors can be treated at once if they cross normally, since the integrability test on the base then separates in normal coordinates.

Let $\pi\colon X\to B$ be a surjective holomorphic submersion of connected compact complex manifolds with connected fibers of dimension $f$, and let $q=\dim B$. Fix a Hermitian form $\omega$ on $X$, a holomorphic line bundle $A$ on $B$ with a smooth Hermitian metric $h$, and $\varphi\in\PSH(X,\omega)$. Write $\theta=c_1(A,h)$, and use holomorphic descent to regard the selected spaces as subspaces of $H^0(B,A^m)$:
\[
 V_m=H^0\bigl(X,\pi^*A^m\otimes\mathcal I(m\varphi)\bigr)
       \subseteq H^0(B,A^m),\qquad m\geq1.
\]
Let $D_1,\ldots,D_N$ be distinct nonempty smooth connected reduced divisors on $B$ whose union has simple normal crossings. For each $i$, choose a smooth metric $H_i$ on $\mathcal O_B(D_i)$, and write $s_i$ for the canonical section and $\eta_i$ for the curvature.

\begin{lemma}\label{lem:anisotropic-fiber-integrability}\index{Multiplier ideal!fiber integrability}
Let $r\geq1$ and $\ell\geq0$ be integers, let $p,b_1,\ldots,b_r$ be positive integers, and let $P>0$. The function
\[
 |z|^{2\ell}\left(|z|^{2p}+\sum_{j=1}^r|w_j|^{2b_j}\right)^{-P}
\]
of $(z,w)\in\mathbb C\times\mathbb C^r$ is locally integrable near the origin if and only if
\[
 \frac{\ell+1}{p}+\sum_{j=1}^r\frac1{b_j}>P.
\]
In the case of equality, the integral diverges logarithmically.
\end{lemma}
\begin{proof}
After integrating over the angles, substitute $x_0=|z|^{2p}$ and $x_j=|w_j|^{2b_j}$. Up to a positive constant, the integral becomes
\[
 \int (x_0+\cdots+x_r)^{-P}x_0^{(\ell+1)/p-1}
                  \prod_{j=1}^r x_j^{1/b_j-1}\,\mathrm dx_0\cdots\mathrm dx_r
\]
over a box near the origin. This box contains a small simplex $\{\sum_jx_j<\delta\}$ and is contained in a larger one. On either simplex use the coordinates $\rho=\sum_jx_j$ and $x_j=\rho y_j$. The integral over $y$ is finite and positive, because all its exponents exceed $-1$, and the integral over $\rho$ is
\[
 \int_0^\delta\rho^{(\ell+1)/p+\sum b_j^{-1}-P-1}\,\mathrm d\rho.
\]
This gives the threshold and the logarithmic divergence at equality. Bounded factors in the weight, holomorphic units and additional tangential variables do not affect this local test.
\end{proof}

\Needspace{36\baselineskip}
\begin{theorem}\label{thm:fiber-integrability-band}\label{thm:crossing-fiber-integrability-band}\label{cor:semipositive-fiber-band}\index{Holomorphic section!fiber integrability band}\index{Holomorphic section!crossing divisorial band}\index{Holomorphic section!semipositive residual}
Suppose that $c_1,\ldots,c_N>0$ and $u\in C^\infty(B,\mathbb R)$ satisfy
\[
 \beta=\theta+\ddc u-\sum_{i=1}^Nc_i\eta_i\geq0,
 \qquad\int_B\beta^q>0,
 \qquad \varphi\geq\sum_{i=1}^Nc_i\log\|\pi^*s_i\|_{H_i}^2-C
                         \quad\hbox{on }X
\]
for a constant $C$. Suppose further that for each $i$ there is one submersion polydisk over an open subset of $D_i$, with coordinates
\[
 \pi(z,t,w,v)=(z,t),\quad D_i=\{z=0\},\quad
 t\in\mathbb C^{q-1},\quad w\in\mathbb C^{r_i},\quad v\in\mathbb C^{f-r_i},
\]
where $1\leq r_i\leq f$, on which
\[
 \varphi\leq\frac{c_i}{p_i}\log\left(|z|^{2p_i}+\sum_{j=1}^{r_i}|w_j|^{2b_{ij}}\right)+C_i
\]
for positive integers $p_i,b_{i1},\ldots,b_{ir_i}$ and a constant $C_i$. Put
\[
 \begin{gathered}
 a_i=p_i\sum_{j=1}^{r_i}b_{ij}^{-1},\quad
 L_{m,i}=\lfloor mc_i\rfloor,\quad
 K_{m,i}=\max\{\lfloor mc_i-a_i\rfloor,0\},\\
 \psi=u+\sum_{i=1}^Nc_i\log\|s_i\|_{H_i}^2.
 \end{gathered}
\]
Then $\psi\in\PSH(B,\theta)$, $\mathrm{MA}_\theta(\psi)=\beta^q$, its mass is $\int_B\beta^q\in(0,\infty)$, and for every $m\geq1$,
\[
 \begin{split}
 W_m&:=H^0\bigl(B,A^m\otimes\mathcal I(m\psi)\bigr)
      =H^0\left(B,A^m\left(-\sum_iL_{m,i}D_i\right)\right)\subseteq V_m,\\
 V_m&\subseteq H^0\left(B,A^m\left(-\sum_iK_{m,i}D_i\right)\right).
 \end{split}
\]
Moreover $\dim V_m-\dim W_m=O(m^{q-1})$. For $q=1$ this difference is at most $\sum_i\lceil a_i\rceil$ for every $m$.
\end{theorem}
\begin{proof}
By the Poincar\'e--Lelong formula, $\theta+\ddc\psi=\beta+\sum_ic_i[D_i]\geq0$. Off the divisors, locality identifies the non-pluripolar top power with $\beta^q$; the non-pluripolar product puts no mass on the divisors, and neither does the smooth measure $\beta^q$. Hence $\mathrm{MA}_\theta(\psi)=\beta^q$, whose mass is finite because $\beta$ is smooth and positive by assumption.

Near a point of $B$, the divisors through that point are distinct coordinate hyperplanes $\{z_i=0\}$, and $\psi=\sum_ic_i\log|z_i|^2+O(1)$, the sum running over these divisors. If a holomorphic germ is divisible by each corresponding $z_i^{L_{m,i}}$, then it is divisible by their product; its residual coefficient is bounded, and since $L_{m,i}+1>mc_i$, Fubini's theorem gives integrability. Conversely, at generic nearby points of each coordinate hyperplane only one divisor passes, and the one-variable integrability condition $\ell+1>mc_i$ forces divisibility by $z_i^{L_{m,i}}$; the divisor of the germ then gives divisibility throughout the chart. This proves $\mathcal I(m\psi)=\mathcal O_B(-\sum_iL_{m,i}D_i)$, also when some $mc_i$ is an integer. Since $u$ is bounded, the global lower bound implies $\varphi\geq\psi\circ\pi-C'$. Hence integrability on the base gives integrability of the pullback in every bounded submersion chart, which proves $W_m\subseteq V_m$.

For the upper inclusion, take a nonzero section in $V_m$ and let $\ell$ be its vanishing order along $D_i$; note that $D_i$ is irreducible, being smooth and connected. In the polydisk over $D_i$ on which the upper bound holds, the coefficient of the section is $z^\ell g(z,t)$, where $g(0,t)$ does not vanish identically on the open subset of $D_i$, since otherwise the identity theorem would increase the order along $D_i$. Shrink the polydisk about a point with $w=0$ at which $g(0,t)\ne0$, so that $g$ becomes a holomorphic unit. The upper bound on $\varphi$ and \cref{lem:anisotropic-fiber-integrability} then force $\ell+1>mc_i-a_i$. For an integer $\ell\geq0$ this is equivalent to $\ell\geq K_{m,i}$. Imposing these orders along all the divisors proves the upper inclusion.

Since $\lfloor mc_i\rfloor-\lfloor mc_i-a_i\rfloor\leq\lceil a_i\rceil$, we have $0\leq L_{m,i}-K_{m,i}\leq\lceil a_i\rceil$. Raise the orders from $(K_{m,i})_i$ to $(L_{m,i})_i$, one divisor and one unit at a time. As in the proof of \cref{thm:horizontal-center-descent}, the restriction sequence along $D_i$ shows that each step loses at most
\[
 h^0\left(D_i,\left.A^m\left(-\sum_j n_jD_j\right)\right|_{D_i}\right),
 \qquad 0\leq n_j\leq mc_j,
\]
dimensions. By \cref{lem:linear-range-section-bound} with $E=A|_{D_i}$ and $F_j=\mathcal O_B(-D_j)|_{D_i}$, this is $O(m^{q-1})$ uniformly over all steps, and there are at most $\sum_i\lceil a_i\rceil$ steps. When $q=1$, the distinct connected divisors are points, and each restriction space has dimension one.
\end{proof}

For a single divisor we write $D,s_D,\eta,c,p,r,b_j,a,L_m,K_m$ for $D_1,s_1,\eta_1,c_1,p_1,r_1,b_{1j},a_1,L_{m,1},K_{m,1}$.

\Needspace{12\baselineskip}
\begin{corollary}\label{cor:exact-anisotropic-fiber-series}
Let $N=1$, and let $c$ and $u$ satisfy the conditions on $\beta$ in \cref{thm:fiber-integrability-band}. Suppose that a nonempty compact smooth submanifold $Z\subset\pi^{-1}(D)$ projects submersively onto $D$, and that near every point of $Z$ there are coordinates as in that theorem in which
\[
 Z=\{z=w_1=\cdots=w_r=0\},\qquad
 \varphi=\frac cp\log\left(|z|^{2p}+\sum_{j=1}^r|w_j|^{2b_j}\right)+O(1),
\]
with fixed $c,p,b_j$. Suppose also that $\varphi$ is locally bounded outside $Z$. Then $\varphi$ satisfies the hypotheses of \cref{thm:fiber-integrability-band}, so its conclusions hold, and moreover
\[
 V_m=H^0(B,A^m(-K_mD))\qquad(m\geq1).
\]
\end{corollary}
\begin{proof}
The local model is bounded below by $c\log|z|^2$ up to a constant. Near points of $\pi^{-1}(D)$ outside $Z$, the function $\varphi$ is locally bounded, so the same lower bound holds there; away from $\pi^{-1}(D)$ both functions are locally bounded. By compactness, finitely many charts give the global lower bound. The upper bound holds in any chart of the model. Conversely, if a section is divisible by $z^{K_m}$, its residual coefficient is bounded and holomorphic in each chart of the model, and the strict inequality $K_m+1>mc-a$ together with \cref{lem:anisotropic-fiber-integrability} gives integrability there. Outside $Z$ the function $\varphi$ is locally bounded, which proves the reverse inclusion.
\end{proof}

Such exact models exist with pole sets of positive dimension inside the fibers. Let $B$ be projective, let $A$ be ample, and let $X=B\times(\mathbb P^1)^f$ with $\pi$ the first projection. Mark a point $o_j$ in each of the first $r$ factors $\mathbb P^1$ of the fiber, and let $E_j\subset X$ be the inverse image of $o_j$. With smooth metrics on the corresponding line bundles, put
\[
 Q=\|\pi^*s_D\|^{2p}+\sum_{j=1}^r\|s_{E_j}\|^{2b_j},
 \qquad \varphi=\lambda\log Q,\qquad c=p\lambda.
\]
The common zero set of the terms of $Q$ is $D\times\{o_1\}\times\cdots\times\{o_r\}\times(\mathbb P^1)^{f-r}$. Near it, the positive factors of the metrics give exactly the local model of \cref{cor:exact-anisotropic-fiber-series}, and $\varphi$ is smooth elsewhere. The inequality $Q\geq\|\pi^*s_D\|^{2p}$ gives the global lower bound. Locally, $\ddc\log Q$ is bounded below by a convex combination of the negatives of the smooth curvature forms of the terms. Hence $\ddc\log Q\geq-C_0\omega$ for a constant $C_0$ off the common zero set, and, after a quadratic correction, the extension theorem for plurisubharmonic functions carries this bound across it. Since $A$ is ample, there is $u$ with $\theta+\ddc u>0$. For sufficiently small $\lambda>0$ we then have $\lambda C_0\leq1$, so that $\varphi\in\PSH(X,\omega)$, and $\theta+\ddc u-p\lambda\eta>0$. Thus the exact series of the corollary also occurs when $r<f$.

\Needspace{14\baselineskip}
\begin{example}\label{ex:anisotropic-fiber-sharp-error}\index{Holomorphic section!anisotropic sharp error}
Let $B=\mathbb P^q$ with a hyperplane $D$, let $A=\mathcal O_B(1)$, and let $X=B\times\mathbb P^1$. Let $\omega_B,\omega_F$ be the Fubini--Study forms normalized to have total mass one, and use the corresponding metrics on the divisor line bundles. For a point $o\in\mathbb P^1$ set
\[
 \varphi=\tfrac14\log\bigl(\|s_D\|^4+\|s_o\|^6\bigr),\qquad
 \omega=\omega_B+\omega_F.
\]
Then $p=2$, $b_1=3$, $c=1/2$, $a=2/3$, and the reference potential $\psi=\tfrac12\log\|s_D\|^2$ has mass $2^{-q}$. The spaces $V_1$ and $W_1$ agree, and for every $m\geq2$ the defect is
\[
 \dim V_m-\dim W_m
    =\binom{\lceil m/2\rceil+q}{q-1}
    \sim\frac{2^{1-q}}{(q-1)!}m^{q-1}.
\]
For $q=1$ the defect is exactly one for every $m\geq2$.
\end{example}
\begin{proof}
Locally, the complex Hessian of the logarithm of a sum of two squared norms of holomorphic sections is bounded below by a convex combination of the negatives of their curvature forms. Hence $\ddc\varphi\geq-\tfrac12\omega_B-\tfrac34\omega_F$ off the pole set, and across it by the extension theorem after a quadratic correction. Thus $\varphi\in\PSH(X,\omega)$. The pole set is $D\times\{o\}$, along which $\varphi$ has the exact local model of \cref{cor:exact-anisotropic-fiber-series}. Here $u=0$, $\beta=\tfrac12\omega_B$, and $K_m=L_m-1$ for $m\geq2$, while $K_1=L_1=0$. The spaces $V_m$ and $W_m$ therefore correspond to homogeneous polynomials of degrees $m-L_m+1$ and $m-L_m$, and since $m-L_m=\lceil m/2\rceil$, the difference of their dimensions is the displayed binomial coefficient.
\end{proof}

In this example $\varphi$ is finite on most of each fiber over $D$ and equals $-\infty$ on a proper subset of it, so $\varphi$ differs from every pullback of a function on $B$ by an unbounded amount. For $p=b_j=1$ and $r=f$, the shift $a=f$ is the shift by the fiber dimension found in \cref{thm:horizontal-center-descent}. \Cref{thm:fiber-integrability-band} thus gives sufficient conditions for the approximation asked for in \cref{qu:fiber-selected-approximation}, and when $B$ is K\"ahler, \cref{prop:descended-kernel-limit} turns them into limits of the kernel measures. We do not know which other fiber singularities admit such an approximation.

\subsection{Several crossing base divisors}
\label{sec:crossing-fiber-integrability}

In \cref{thm:fiber-integrability-band} the upper bounds are required only on one polydisk over each divisor, which can be chosen away from the crossings; integrability at the crossings is controlled by the global lower bound. Exact models can nevertheless have fiber poles over the crossings, as in the following example, where different divisors use separate fiber coordinates.

\Needspace{32\baselineskip}
\begin{example}\label{ex:crossing-fiber-sharp-error}\index{Holomorphic section!crossing sharp error}
Let $B=(\mathbb P^1)^q$, $X=B\times(\mathbb P^1)^q$ and $A=\mathcal O_B(1,\ldots,1)$. For each $i$, let $D_i\subset B$ be the inverse image of a marked point in the $i$th factor of $B$, and let $E_i\subset X$ be the inverse image of a marked point in the $i$th factor of the fiber. Use normalized Fubini--Study metrics on the corresponding line bundles, and write $\omega_{B,i},\omega_{F,i}$ for the pullbacks to $X$ of the normalized Fubini--Study forms on the $i$th factors of $B$ and of the fiber. Put
\[
 \begin{gathered}
 R_i=\|\pi^*s_i\|^2,\quad T_i=\|s_{E_i}\|^2,\quad
 \varphi=\frac14\sum_{i=1}^q\log(R_i^2+T_i^3),\\
 \omega=\sum_i(\omega_{B,i}+\omega_{F,i}),\qquad
 \psi=\frac12\sum_i\log\|s_i\|^2.
 \end{gathered}
\]
Then $\varphi\in\PSH(X,\omega)$, and \cref{thm:fiber-integrability-band} applies with $N=q$, $u=0$, $c_i=1/2$, $p_i=2$, $b_{i1}=3$ and $a_i=2/3$. The reference potential $\psi$ has mass $q!/2^q$. With $L_m=\lfloor m/2\rfloor$ and $K_m=\max\{\lfloor m/2-2/3\rfloor,0\}$, the selected and reference spaces are exactly
\[
 \begin{split}
 V_m&=H^0\bigl(B,\mathcal O_B(m-K_m,\ldots,m-K_m)\bigr),\\
 W_m&=H^0\bigl(B,\mathcal O_B(m-L_m,\ldots,m-L_m)\bigr).
 \end{split}
\]
They agree for $m=1$, and for $m\geq2$ the defect is
\[
 \dim V_m-\dim W_m
       =(\lceil m/2\rceil+2)^q-(\lceil m/2\rceil+1)^q.
\]
For $q\geq2$ this is $q2^{1-q}m^{q-1}+O(m^{q-2})$, and for $q=1$ it is exactly one.
\end{example}
\begin{proof}
By the calculation in \cref{ex:anisotropic-fiber-sharp-error}, the complex Hessian of each summand of $\varphi$ is bounded below by $-\tfrac12\omega_{B,i}-\tfrac34\omega_{F,i}$, also across its pole set by the extension theorem after a quadratic correction. Summing proves $\varphi\in\PSH(X,\omega)$. The global lower bound follows from $R_i^2+T_i^3\geq R_i^2$. For the upper bound over $D_i$, away from the other divisors, place the $i$th fiber coordinate at its marked point; all the other summands are smooth there. The residual is $\beta=\tfrac12\sum_i\omega_{B,i}>0$, whose top integral is $q!/2^q$.

The theorem gives the required generic orders along every $D_i$. Conversely, near a point where several of the pairs $(R_i,T_i)$ vanish, the density $\mathrm{e}^{-m\varphi}$ is comparable to the product of the factors $(|z_i|^4+|w_i|^6)^{-m/4}$ in disjoint pairs of variables. A descended coefficient divisible by every $z_i^{K_m}$ has a bounded holomorphic residual. Each factor is integrable because $K_m+1>m/2-2/3$, and Fubini's theorem gives integrability at the intersections; elsewhere the same argument uses fewer factors. This proves the identity for $V_m$, and the theorem gives $W_m$. A line bundle $\mathcal O_B(d,\ldots,d)$ with $d\geq0$ has $(d+1)^q$ independent sections, the products of monomials in the separate factors. Since $K_m=L_m-1$ for $m\geq2$ and $m-L_m=\lceil m/2\rceil$, subtraction gives the displayed defect and its leading term.
\end{proof}

As before, the selecting potential takes finite values and the value $-\infty$ in the fibers over the divisors, so it does not descend.

\subsection{Fiber integrability on a modified base}
\label{sec:modified-base-fiber-integrability}

Retain $\pi:X\to B$, $(A,h)$, $\theta$, $\varphi$ and $V_m$ from \cref{sec:fiber-integrability-band}. A modification of the base can turn a higher-codimension singularity into divisorial data. The integrability test then acts on the pulled-back section multiplied by the Jacobian. This fixed factor changes the exact orders, while the restriction estimate still bounds the section loss.

Let $\rho:Y\to B$ be a modification of connected compact complex manifolds. Form the fiber product and its projections
\[
 X'=X\times_B Y,\qquad \pi':X'\to Y,\qquad \mu:X'\to X.
\]
Local submersion coordinates show that $X'$ is smooth, $\pi'$ is a submersion with the same compact connected fibers, and $\mu$ is a modification. The relative canonical bundle $R_\rho$ and Jacobian section $J_\rho$ are those of \cref{sec:psh-birational-ideals}. The coordinate Jacobian of $\mu$ is the base Jacobian times the identity in the fiber variables, giving
\[
 R_\mu=\pi'^*R_\rho,\qquad J_\mu=\pi'^*J_\rho.
\]
Put $\varphi'=\varphi\circ\mu$ and $\theta'=\rho^*\theta$. Choose distinct nonempty smooth connected reduced divisors $D_1,\ldots,D_N$ on $Y$ with simple-normal-crossing union, smooth metrics $H_i$, defining sections $s_i$ and curvatures $\eta_i$.

\Needspace{14\baselineskip}
\begin{theorem}\label{thm:modified-base-fiber-band}\index{Holomorphic section!modified base band}
Suppose $c_i>0$ and $u\in C^\infty(Y,\mathbb R)$ satisfy
\[
 \beta=\theta'+\ddc u-\sum_i c_i\eta_i>0,
 \qquad \varphi'\geq\sum_i c_i\log\|\pi'^*s_i\|_{H_i}^2-C.
\]
For each $D_i$, assume one anisotropic upper witness for $\varphi'$ and $\pi'$ away from the other divisors, as in \cref{thm:crossing-fiber-integrability-band}, with fixed $p_i,b_{ij}$. Define $a_i,L_{m,i},K_{m,i}$ by the formulas in that theorem, and put
\[
 \psi'=u+\sum_i c_i\log\|s_i\|_{H_i}^2,
 \qquad M_m=\rho^*A^m\otimes R_\rho.
\]
There is a unique $\psi\in\PSH(B,\theta)$ with $\psi\circ\rho=\psi'$, and
\[
 \mathrm{MA}_\theta(\psi)=\rho_*(\beta^q),\qquad
 \int_B\mathrm{MA}_\theta(\psi)=\int_Y\beta^q\in(0,\infty).
\]
The map $T_m(s)=J_\rho\rho^*s$ identifies $H^0(B,A^m)$ with $H^0(Y,M_m)$. For $W_m=H^0(B,A^m\otimes\mathcal I(m\psi))$,
\[
 \begin{split}
 T_m(W_m)&=H^0\left(Y,M_m\left(-\sum_iL_{m,i}D_i\right)\right),\\
 T_m(W_m)&\subseteq T_m(V_m)
       \subseteq H^0\left(Y,M_m\left(-\sum_iK_{m,i}D_i\right)\right).
 \end{split}
\]
Consequently $W_m\subseteq V_m$ and $\dim V_m-\dim W_m=O(m^{q-1})$. On a curve the defect is at most $\sum_i\lceil a_i\rceil$ for every $m$.
\end{theorem}
\begin{proof}
Poincar\'e--Lelong gives $\theta'+\ddc\psi'=\beta+\sum_i c_i[D_i]$. Thus $\psi'\in\PSH(Y,\rho^*\theta)$, and \cref{prop:psh-descent-modification} gives the unique descended potential. The non-pluripolar measure upstairs is $\beta^q$ by locality off the divisor union and exclusion of that union. Apply the measure identity in \cref{prop:birational-partial-envelopes} to obtain the asserted pushforward and mass.

\Cref{lem:adjoint-ideal-transform}, first with weight zero and then with $m\psi$, gives the whole-space identification and the identity for $T_m(W_m)$. The simple-normal-crossing integrability calculation in \cref{thm:crossing-fiber-integrability-band} supplies $\mathcal I(m\psi')=\mathcal O_Y(-\sum_iL_{m,i}D_i)$. Boundedness of $u$ and the global lower bound give $\varphi'\geq\psi'\circ\pi'-C'$. Since $\rho\circ\pi'=\pi\circ\mu$ and $\mu$ is surjective, this descends to $\varphi\geq\psi\circ\pi-C'$. Integrability in bounded submersion charts proves $W_m\subseteq V_m$.

For $s\in V_m$, the adjoint transformation for $\mu$ gives
\[
 J_\mu\mu^*s=\pi'^*T_m(s)
 \in H^0\bigl(X',\pi'^*M_m\otimes\mathcal I(m\varphi')\bigr).
\]
Its descended coefficient in each upper witness therefore has order at least $K_{m,i}$, by \cref{lem:anisotropic-fiber-integrability} and the generic unit-coefficient argument of \cref{thm:fiber-integrability-band}. This order is the order of $T_m(s)$, including its Jacobian factor. It proves the upper inclusion.

The divisor restriction chain on $Y$ has at most $\sum_i\lceil a_i\rceil$ steps. On $D_i$, each target twists $(\rho^*A|_{D_i})^m$ by the powers $\mathcal O_Y(-D_j)|_{D_i}$ in a linear range and by the fixed bundle $R_\rho|_{D_i}$. Include this last bundle with exponent one in \cref{lem:finite-family-section-bound}; every target has dimension $O(m^{q-1})$ uniformly. Summing gives the estimate. For $q=1$, each divisor is a point and every target has dimension one.
\end{proof}

The theorem uses a specified smooth modification and divisorial data with a positive residual. It does not assume that arbitrary selecting potentials furnish these data. The following point-pole model makes both the residual and the Jacobian correction explicit.

\Needspace{14\baselineskip}
\begin{example}\label{ex:modified-point-fiber-series}\index{Holomorphic section!point pole Jacobian}
Let $q\geq2$, $B=\mathbb P^q$, $A=\mathcal O_B(1)$, $p=[1:0:\cdots:0]$ and $X=B\times\mathbb P^1$. In normalized Fubini--Study metrics, for a marked fiber point $o$ put
\[
 R=\frac{\sum_{j=1}^q|Z_j|^2}{\sum_{j=0}^q|Z_j|^2},\quad
 T=\|s_o\|^2,\quad \varphi=\tfrac14\log(R^2+T^3),\quad
 \psi=\tfrac12\log R.
\]
Then $\varphi\in\PSH(X,\omega_B+\omega_F)$ and $\psi\in\PSH(B,\omega_B)$ has mass $1-2^{-q}$. With the ideal sheaf $\mathfrak m_p$ of $p$, the exact spaces are
\[
 \begin{gathered}
 V_m=H^0\bigl(B,A^m\otimes\mathfrak m_p^{k_m}\bigr),\qquad
 W_m=H^0\bigl(B,A^m\otimes\mathfrak m_p^{\ell_m}\bigr),\\
 k_m=\max\{\lfloor m/2-2/3\rfloor-q+1,0\},\qquad
 \ell_m=\max\{\lfloor m/2\rfloor-q+1,0\}.
 \end{gathered}
\]
Their defect vanishes for $m<2q$. For $m\geq2q$ it is
\[
 \dim V_m-\dim W_m=\binom{\lfloor m/2\rfloor-1}{q-1}
       \sim\frac{2^{1-q}}{(q-1)!}m^{q-1}.
\]
\end{example}
\begin{proof}
A common denominator writes $R^2+T^3$ as a sum of squared holomorphic section coefficients in $\mathcal O_X(2,3)$ divided by its smooth metric denominator. Hence $\ddc\varphi\geq-\tfrac12\omega_B-\tfrac34\omega_F$, including across the common zero set by local psh extension.

Let $\rho:Y\to B$ be the blowup of $p$, with exceptional divisor $E$, and let $\tau:Y\to\mathbb P^{q-1}$ be the resolved projection $[Z_0:\cdots:Z_q]\mapsto[Z_1:\cdots:Z_q]$. The graph embedding $(\rho,\tau)$ shows that
\[
 \alpha=\rho^*\omega_B,\quad\gamma=\tau^*\omega_{\mathbb P^{q-1}},\qquad
 \beta=\tfrac12(\alpha+\gamma)>0.
\]
The divisor bundle is $\mathcal O_Y(E)=\rho^*\mathcal O_B(1)\otimes\tau^*\mathcal O_{\mathbb P^{q-1}}(-1)$. Its quotient metric has curvature $\alpha-\gamma$ and satisfies $R\circ\rho=\|s_E\|^2$. Thus the theorem applies with $u=0$, $c=1/2$, $a=2/3$, and residual $\beta$. The pulled-back selector has the exact model $\tfrac14\log(|z|^4+|w|^6)+O(1)$ along $E\times\{o\}$ and is locally bounded elsewhere. Divisibility by $z^{K_m}$, where $K_m=\max\{\lfloor m/2-2/3\rfloor,0\}$, gives the reverse upper inclusion by \cref{lem:anisotropic-fiber-integrability} in every pole chart.

In blowup coordinates $Z_1/Z_0=z$, $Z_j/Z_0=zt_j$ for $j\geq2$, the Jacobian is $z^{q-1}$. A section of base order $n$ therefore gives $T_m(s)$ of generic order $n+q-1$ along $E$. The two divisorial thresholds yield exactly $k_m,\ell_m$ above. Homogeneous monomials realize every jet of total degree at most $m$ at $p$. For $m\geq2q$, $k_m=\ell_m-1$ and the newly allowed homogeneous degree $\ell_m-1$ has $\binom{\ell_m+q-2}{q-1}$ monomials, giving the defect. For $m<2q$ both orders are zero.

Finally, in cohomology $[\beta]=\rho^*c_1(A)-\tfrac12[E]$. The mixed intersection numbers containing both $\rho^*c_1(A)$ and $[E]$ vanish because the pullback restricts to zero on $E$. The normal bundle of $E\simeq\mathbb P^{q-1}$ is $\mathcal O_E(-1)$, so $\int_Y[E]^q=(-1)^{q-1}$, while $\int_Y(\rho^*c_1(A))^q=1$. Expanding gives $\int_Y\beta^q=1-2^{-q}$, which is the mass of $\psi$ by the theorem.
\end{proof}

The selector is finite at most points of the fiber over $p$ and equals $-\infty$ at $o$, so it does not descend up to bounded difference. Its base reference has a higher-codimension pole; the Jacobian is what converts the resolved divisor orders into the base point-ideal orders.



